% Generated draft: requires mathematical and source review.
\documentclass[11pt]{article}
\usepackage[margin=25mm]{geometry}
\usepackage{fontspec}
\setmainfont{Latin Modern Roman}
\usepackage{amsmath,amssymb,mathtools,unicode-math}
\setmathfont{Latin Modern Math}
\usepackage{enumitem,needspace,xurl,hyperref}
\setlength{\parindent}{0pt}
\setlength{\parskip}{5pt}
\setlength{\emergencystretch}{3em}
\newcommand{\N}{\mathbb{N}}
\newcommand{\Z}{\mathbb{Z}}
\newcommand{\Q}{\mathbb{Q}}
\newcommand{\R}{\mathbb{R}}
\newcommand{\C}{\mathbb{C}}
\newcommand{\F}{\mathbb{F}}
\newcommand{\A}{\mathbb{A}}
\newcommand{\PP}{\mathbb{P}}
\newcommand{\E}{\mathbb{E}}
\newcommand{\eps}{\varepsilon}
\newcommand{\dd}{\,\mathrm{d}}
\begin{document}
\section*{UnsolvedMath ID: 800183}
\textbf{Source problem code:} catalogue entry
\subsection{Definitions and mathematical statement}
Using the standard ternary divisor-function convention, define, for each positive integer \(n\),
\[
d_3(n):=\#\{(a,b,c)\in\Z_{>0}^{3}:abc=n\},
\]
where triples are ordered. For real \(x\geq 1\), let
\[
S(x):=\sum_{\substack{n\in\Z_{>0}\\n\leq x}}d_3(n)d_3(n+1).
\]
Determine an explicit leading asymptotic formula for \(S(x)\) as \(x\to+\infty\). The supplied statement does not define \(d_3\); the standard convention above is an explicit interpretive assumption. It also does not specify a proposed main term, whether lower-order terms are required, or a desired error bound.
\subsection{Short English statement}
Determine the asymptotic growth of the sum of products of the ternary divisor function at consecutive positive integers. The source labels the problem folklore and records its status as open; that status is supplied metadata, not an independently verified assessment.
\subsection{Sources}
\begin{itemize}
\item Supplied local UnsolvedMath registry record: permanent ID \(800183\), source problem code \texttt{catalogue entry}, in the AMR Open Problem Lists collection. The record attributes the problem to Hamieh, Kadiri, and Ng, with proposed year 2024, and labels it open.
\item Hamieh, Kadiri, and Ng, \emph{Moments of L-functions Problem List} (2024), Problem 35, under Moments of the Riemann zeta function: \url{https://arxiv.org/abs/2405.17800}. The supplied statement includes the attribution ``folklore.''
\item Supplied provenance: extraction method \texttt{arxiv\_latex}; source access date \texttt{2026-10-16}. This access date is supplied provenance; no independent source visit or literature verification is asserted here.
\end{itemize}
\end{document}
